
\linksection{fixtures}{Fixtures}

pytest calls the fixture function to inject a dependency into the test function (with an arg of the same name):
\vspace{-.5em}

\begin{minted}[autogobble,escapeinside=||]{python}
    @pytest.fixture
    def |\highlightbox[highlightcolor!80!black]{rpn}|() -> RPNCalc:
        """RPN calculator with a default config."""
        return RPNCalc(Config())

    def test_empty_stack(|\highlightbox{rpn}|: RPNCalc):
        assert rpn.stack == []
\end{minted}

Imagine pytest running your test like: \\[.25em]
{\color{codeblue}\mintinline[escapeinside=||]{python}{test_empty_stack(|\highlightbox{rpn\vphantom{()}}|=|\highlightbox[highlightcolor!80!black]{rpn()}|)}}
\vspace{.5em}

Useful pytest arguments:
\vspace{-.5em}

\begin{description}[itemsep=0pt]
    \item[\mono{--fixtures}] List all available fixtures
    \item[\mono{--fixtures-per-test}] \strut\\ Show which fixtures are used by which tests
    \item[\mono{--setup-show}] \strut\\ Show setup/teardown of fixtures during test
\end{description}

\linksubsection{conftest}{conftest.py}

You can move fixture functions into \mono{conftest.py}:

\begin{itemize}[itemsep=0pt]
    \item Automatically discovered by pytest.
    \item Visible to modules in same or sub directory.
    \item Put hooks and cross-module used fixtures into conftest file(s) but \textbf{don't import} from them \\ (exception: classes for type annotations).
\end{itemize}

\dirtree{%
    .1 myproject/tests/.
    .2 export/.
    .3 \highlightbox{conftest.py}\vspace{-1pt}.
    .3 \highlightbox{test\_jsonexport.py}.
    .3 \highlightbox[white]{test\_yamlexport.py}.
    .2 \highlightbox[white]{conftest.py}.
    .2 \highlightbox[highlightcolorred]{test\_utils.py}.
    .2 \codeskip*.
}
